%-----------------------------------------------------------------------
\section{Verification}
\label{sec:verify}
%-----------------------------------------------------------------------

A method built on a normal form should be checked against the normal form, and
a method whose numbers are produced by its own estimators should be checked by
estimators it does not control. This section does both, and then reports the
two errors in our own earlier work that the checks exposed.

\subsection{The profile is the true Airy function}
\label{sec:verify-profile}

The implementation evaluates the fold profile through a 2048-entry table. We
recomputed the table's defining integral independently and compared it against
\verb|scipy.special.airy|. The maximum absolute deviation from $\Ai(-u)$,
normalised to the peak, is \num{8.118e-03}, so the shipped table is accurate to
under one percent of peak. The first sign change agrees exactly between the
table and the independent evaluation, at $u = 2.2915$, against the analytic first
zero of $\Ai(-u)$ at $u = 2.3381$; the residual is the table's grid spacing at
that point, not an error in the profile.

\subsection{The fold pattern peaks where $|\Ai|^2$ says}
\label{sec:verify-peak}

The profile's maximum in the solver sits at $X = 1.00$ where the analytic
maximum of $|\Ai(-X)|^2$ is at $X = 1.0188$ (Table~\ref{tab:airyfacts}). The
agreement is to the resolution at which the peak position was measured, and it
is a genuine check rather than a tautology: a profile built from the wrong
canonical scaling would peak at a different $X$, and the fact that ours peaks
beside the fold rather than on it is the qualitative difference from the
geometric model set out in Section~\ref{sec:foldnormal}.

\subsection{The $\lambda^{2/3}$ law}
\label{sec:verify-scaling}

Corollary~\ref{cor:scaling} predicts an exponent of $2/3 = 0.6667$ in
\eqref{eq:scalelaw}. Fitting the exponent from the solver's own per-channel
widths gives \num{0.6977}. This is the load-bearing number for the dispersion
story, so it is worth being precise about what the discrepancy means: the fitted
value comes from three channels spanning a factor of only \num{1.44} in
wavelength, which is a short baseline for an exponent fit, and the fit's
uncertainty at that baseline is comparable to the deviation. The derivation,
not the fit, is what fixes the exponent; the fit is a consistency check that
passes within its resolution.

\subsection{The ray map survives a second implementation}
\label{sec:verify-parity}

The C++ baker and the Python reference were written separately from the same
definitions and compared at $512^2$ with \texttt{airy\_scale}${}=1.0$:

\begin{table}[t]
\centering
\caption{C++ against Python reference at $512^2$. The wave-channel differences
are localised on fold pixels and come from the table resolution in C++ against
exact \texttt{scipy.special.airy}, plus Euclidean-distance-transform
tie-breaking between equidistant band pixels.}
\label{tab:parity}
\small
\begin{tabular}{@{}lrrr@{}}
\toprule
Field & C++ peak & Python peak & Max abs diff \\
\midrule
geometric & 9.281 & 9.280 & \num{5.6e-04} \\%  (0.006%)
wave R    & 4.754 & 4.731 & \num{6.1e-01} \\
wave G    & 4.762 & 4.749 & \num{6.1e-01} \\
wave B    & 4.774 & 4.772 & \num{6.1e-01} \\
\bottomrule
\end{tabular}
\end{table}

Peaks agree to \SI{0.3}{\percent}. Parallel against serial output agrees to
\num{3.8e-06} absolute (\num{1.1e-08} relative), which is floating-point
reduction order and nothing else. Two independent implementations of the ray map
agreeing to this level is the check that the map \eqref{eq:raymap} is
implemented as specified rather than as remembered.

\subsection{Energy audit: the non-circular statement}
\label{sec:verify-energy}

Remark~\ref{rem:circular} forbids using the ratio $k^{*}$ as evidence about
energy. The non-circular measurement follows the light path through every
interface with the Fresnel weights of \eqref{eq:fresnel} and accounts for where
the power went. The result is a transport throughput of \SI{96.41}{\percent} at
the three wavelengths, with Fresnel reflection as the only loss channel: no
energy is created or destroyed by the correction, and the \SI{3.59}{\percent} of
loss is accounted for by the interfaces.

The same audit measures what the fold treatment does to the field's energy
rather than to the image's appearance: the wave-corrected field retains on the
order of \SI{28}{\percent} of the geometric field's energy, the remainder being
removed by the ceiling and by the profile's own boundedness. This is the honest
global statement the method can support, and it is a statement about a bounded
display field rather than about transport.

\subsection{An index model that was wrong}
\label{sec:verify-index}

An earlier version of this work used a fitted Cauchy dispersion model instead of
the measured data of Section~\ref{sec:dispersion}. Evaluated at the render
channels, that fit carried roughly \num{1.97} times too much dispersion in the
water and gave a glass index about \num{0.0017} low at one channel. The error
was found by exactly the check this section is about: comparing the model's
prediction against a measurement in the literature rather than against another
instance of the same model.

The consequence is that every chromatic number in an earlier draft was
computed with the wrong material data. The numbers in this paper are not: all
of them use \eqref{eq:sellmeier} and \eqref{eq:rubin} with the coefficients of
Daimon and Masumura \cite{daimon2007} and Rubin \cite{rubin1985}. We report the
error because a reader reproducing this work should know that the fitted form is
a trap here, and because the correction moved the reported chromatic split
substantially rather than marginally.

\subsection{A metric that manufactured a false refutation}
\label{sec:verify-metric}

The most instructive error in this project was not physical. When SPPM was first
scored against the ground truth it returned $r = -0.046$, which reads as a
refutation of SPPM. It was an artefact of two decisions, both in the metric.

\textbf{Sign convention.} SPPM's caustic is isolated as
$\text{optics} - \text{no-optics}$, and that differential is not non-negative:
removing the glass also removes its shadow, so \SI{11.4}{\percent} of pixels go
negative with $|\text{neg}|/|\text{pos}| = 1.40$. Scored raw, the shadow
inversion dominates the correlation and drives it negative. Scored on its
positive part, which is the caustic, SPPM reads $+0.853$. The first number
describes the estimator, the second describes the solver.

\textbf{Tile naming.} In these tools the tile argument is a tile size in
pixels, while the tables in our evidence file name the \SI{40}{\milli\meter} row
by its tile count per side. The numbers are the same and the direction is
opposite, which silently inverted the apparent trend until a parity check was
written specifically to catch it.

A claim that flips with the estimator convention is not a result. Both
conventions are printed by the tool
(\texttt{tools/metric\_parity\_check.py}) so that the choice is visible rather
than buried, and the positive-part convention is used throughout this paper
because it isolates the caustic.
